\documentclass[11pt]{article}

\usepackage[margin=1in]{geometry}
\usepackage{amsmath,amssymb,amsfonts}
\usepackage[hidelinks]{hyperref}

\newcommand{\Spin}{\mathrm{Spin}}
\newcommand{\U}{\mathrm{U}}
\newcommand{\cO}{\mathcal{O}}

\title{Route D2-E3: E3-Instanton and Green--Schwarz Audit for \(\zeta K_{\rm tr}\)}
\author{BoYu Chen}
\date{June 11, 2026}

\begin{document}
\maketitle

\section*{Purpose}

This note tests the second Route-D string-constrained graft: whether the
Route-B Majorana contact
\begin{equation}
  \frac{M_R}{M_*}=M_V+\zeta K_{\rm tr}
\end{equation}
can be interpreted as an Euclidean D3-brane instanton effect in a Type IIB/F
theory compactification.  The goal is not to derive the benchmark value of
\(\zeta\) to many digits.  The goal is to check whether the structure
\(\zeta K_{\rm tr}\), the hidden phase/radial language, and the
Majorana-only selection rule have a standard string-theory interpretation.

The answer is conditionally yes: an E3 instanton can supply a holomorphic
factor
\begin{equation}
  \zeta_\Gamma=A_\Gamma \exp(2\pi i T_\Gamma),
  \qquad
  T_\Gamma=\theta_\Gamma+i\,t_\Gamma,
  \label{eq:e3-factor}
\end{equation}
with modulus controlled by the volume-like variable \(t_\Gamma\) and phase
controlled by the axion-like variable \(\theta_\Gamma\).  Charged zero modes
and Green--Schwarz/Stueckelberg selection rules can allow a perturbatively
forbidden Majorana operator while forbidding other perturbative operators.
However, the existence of the required instanton and its zero-mode spectrum is
global compactification data.

\section{Instanton coefficient as hidden phase/radial data}

In Type IIB/F-theory language, an E3 instanton wrapping a suitable divisor or
cycle \(\Gamma\) contributes to the superpotential with a holomorphic factor
of the schematic form
\begin{equation}
  A_\Gamma\,e^{-S_\Gamma+i\varphi_\Gamma}.
\end{equation}
Equivalently, with the complexified modulus convention
\begin{equation}
  T_\Gamma=\theta_\Gamma+i\,t_\Gamma,
\end{equation}
one may write Eq.~\eqref{eq:e3-factor}.  The two descriptions are related by
\begin{equation}
  S_\Gamma-\log|A_\Gamma|=2\pi t_\Gamma,\qquad
  \varphi_\Gamma=\arg A_\Gamma+2\pi\theta_\Gamma .
\end{equation}
Thus the Route-B phrase ``hidden phase/radial sector'' maps directly to an
axion-volume pair:
\begin{equation}
  \text{phase}\leftrightarrow \theta_\Gamma,\qquad
  \text{radial suppression}\leftrightarrow t_\Gamma.
\end{equation}

For the current benchmark,
\begin{equation}
  \zeta_{\rm tar}=0.1076472949+0.0736514853\,i,
\end{equation}
so
\begin{align}
  |\zeta_{\rm tar}|&=0.13043190325293763,\\
  \arg\zeta_{\rm tar}&=0.600038020318215.
\end{align}
If the instanton prefactor has unit modulus, the effective action needed to
produce the modulus is
\begin{equation}
  S_{\rm eff}=-\log|\zeta_{\rm tar}|
  =
  2.0369040025655094.
  \label{eq:seff}
\end{equation}
In the convention \(\zeta=\exp(2\pi iT_\Gamma)\), this corresponds to
\begin{equation}
  \operatorname{Im}T_\Gamma
  =
  \frac{S_{\rm eff}}{2\pi}
  =
  0.3241833406119675
\end{equation}
if \(|A_\Gamma|=1\).  These are moderate nonperturbative data; they should not
be confused with a small discrete clockwork prediction.

\section{Majorana-only selection rule}

Perturbative D-brane/F-theory vacua often contain abelian symmetries that act
as selection rules on the superpotential.  Some of these \(\U(1)\) symmetries
can become massive through a Green--Schwarz/Stueckelberg mechanism, yet remain
as perturbative selection rules.  An instanton factor can transform under the
same \(\U(1)\) because its axion shifts.  This allows a perturbatively
forbidden charged operator to appear nonperturbatively if the total charge of
the instanton factor and the inserted fields is zero.

For the right-handed-neutrino Majorana operator, a minimal charge bookkeeping
model is
\begin{equation}
  q(N)=+1,\qquad q(NN)=+2,\qquad q(e^{-S_\Gamma+i\varphi_\Gamma})=-2.
\end{equation}
Then
\begin{equation}
  q\left(e^{-S_\Gamma+i\varphi_\Gamma}NN\right)=0,
\end{equation}
while \(NN\) alone is perturbatively forbidden.  In a microscopic instanton
calculation, this charge bookkeeping is implemented by charged fermionic zero
modes at the intersections of the instanton with the seven-brane stack.  The
zero modes must be soaked up by insertions of the desired fields.  If the only
available charged-zero-mode saturation pattern inserts \(N_iN_j\), the
instanton generates a Majorana mass but not a Dirac Yukawa or colored-triplet
operator.

This is the string-theory version of the Route-B Majorana-only matching rule.
It is stronger than an arbitrary EFT prohibition, but it remains conditional:
the divisor \(\Gamma\), its orientifold/F-theory lift, its universal zero
modes, and its charged zero modes must be checked in a concrete
compactification.

\section{Direction in family space}

Route B separates the tensor direction from the complex coefficient:
\begin{equation}
  \Delta M_R/M_*=\zeta K_{\rm tr}.
\end{equation}
The E3 instanton naturally explains the coefficient as
\begin{equation}
  \zeta=\zeta_\Gamma=A_\Gamma\exp(2\pi iT_\Gamma).
\end{equation}
The tensor direction \(K_{\rm tr}\) requires an additional locality or
equivariance assumption.  If the instanton data are blind to the detailed
coordinate on the family curve and respect the same local \(SL(2)\)-covariant
structure used in the Route-A carrier
\begin{equation}
  V=H^0(\mathbb{CP}^1,\cO(2))\simeq \mathrm{Sym}^2\mathbb C^2,
\end{equation}
then the only invariant symmetric tensor direction in the contact sector is
the second transvectant:
\begin{equation}
  K_{\rm tr}\in\mathrm{Sym}^0
  \subset
  \mathrm{Sym}^2(\mathrm{Sym}^2\mathbb C^2)
  =
  \mathrm{Sym}^4\mathbb C^2\oplus\mathrm{Sym}^0.
\end{equation}
Thus a family-blind instanton can consistently land on \(K_{\rm tr}\).  This
is not a proof from exact global \(SL(2)\) flavor symmetry.  Quantum gravity
does not allow exact ungauged global symmetries.  The correct statement is
local and spurionic: if the instanton is insensitive to the Veronese coordinate
data except through the contact pairing, the unique available invariant
contact direction is \(K_{\rm tr}\).

\section{Why no ten-digit prediction follows}

The instanton interpretation naturally explains the form, selection rule, and
order of magnitude of \(\zeta\), but it does not compute the benchmark complex
number to high precision.  Multi-instanton or multi-cover corrections are
expected at relative scales set by powers of the one-instanton suppression.
For the benchmark,
\begin{equation}
  |\zeta|^2=0.01701248138618368.
\end{equation}
The contact-sensitivity window recorded in the Route-B benchmark scan was
\begin{equation}
  \Delta s=3.208798\times10^{-5}.
\end{equation}
The ratio is
\begin{equation}
  \frac{|\zeta|^2}{\Delta s}=530.1823731560442.
\end{equation}
Therefore, unless a concrete compactification shows that the prefactor and
higher instanton corrections are exceptionally controlled, the E3 picture
cannot justify a ten-digit prediction of \(\zeta\).  It can justify the
conditional structure
\begin{equation}
  \zeta K_{\rm tr}
  =
  A_\Gamma e^{2\pi iT_\Gamma}K_{\rm tr}
\end{equation}
and a moderate effective action \(S_{\rm eff}\simeq2.04\).

\section{Result}

The D2-E3 graft is viable as a conditional string-origin statement:
\begin{align}
  \text{E3 instanton on }\Gamma
  &\leadsto
  \zeta_\Gamma=A_\Gamma e^{2\pi iT_\Gamma},\\
  T_\Gamma=\theta_\Gamma+i t_\Gamma
  &\leadsto
  \text{hidden phase/radial data},\\
  \text{GS/Stueckelberg charge shift plus charged zero modes}
  &\leadsto
  \text{possible Majorana-only selection},\\
  \text{family-blind contact locality}
  &\leadsto
  K_{\rm tr}\text{ as the unique contact direction}.
\end{align}
It is not a completed global model.  To promote this from an appendix-level
interpretation to a UV completion, one must construct the instanton divisor,
count universal and charged zero modes, check that unwanted operators are
absent, and replay the Route-B canonical-normalization audit.

\begin{thebibliography}{9}

\bibitem{Witten1996}
E.~Witten,
``Non-perturbative superpotentials in string theory,''
\href{https://arxiv.org/abs/hep-th/9604030}{arXiv:hep-th/9604030}.

\bibitem{BlumenhagenCveticWeigand2006}
R.~Blumenhagen, M.~Cvetic, and T.~Weigand,
``Spacetime instanton corrections in 4D string vacua: The seesaw mechanism
for D-brane models,''
\href{https://arxiv.org/abs/hep-th/0609191}{arXiv:hep-th/0609191}.

\bibitem{IbanezUranga2006}
L.~E.~Ibanez and A.~M.~Uranga,
``Neutrino Majorana masses from string theory instanton effects,''
\href{https://arxiv.org/abs/hep-th/0609213}{arXiv:hep-th/0609213}.

\bibitem{CveticRichterWeigand2007}
M.~Cvetic, R.~Richter, and T.~Weigand,
``Computation of D-brane instanton induced superpotential couplings:
Majorana masses from string theory,''
\href{https://arxiv.org/abs/hep-th/0703028}{arXiv:hep-th/0703028}.

\bibitem{BianchiCollinucciMartucci2011}
M.~Bianchi, A.~Collinucci, and L.~Martucci,
``Magnetized E3-brane instantons in F-theory,''
\href{https://arxiv.org/abs/1107.3732}{arXiv:1107.3732}.

\bibitem{CveticDonagiHalversonMarsano2012}
M.~Cvetic, R.~Donagi, J.~Halverson, and J.~Marsano,
``On seven-brane dependent instanton prefactors in F-theory,''
\href{https://arxiv.org/abs/1209.4906}{arXiv:1209.4906}.

\end{thebibliography}

\end{document}
